\documentclass[../../../include/open-logic-section]{subfiles}

\begin{document}

\olfileid[pt]{sth}{spine}{rank}
\olsection{Posto}

Tendo definido os estágios como os $V_\alpha$ e sabendo
que todo o conjunto é subconjunto de algum estágio, podemos definir o
\emph{posto} de um conjunto. Intuitivamente, o posto de $A$ é o primeiro momento
em que $A$ é formado. Com mais precisão:

\begin{defn}\ollabel{defnsetrank}
Para cada conjunto $A$, $\setrank{A}$ é o menor ordinal $\alpha$ tal que $A
\subseteq V_\alpha$.
\end{defn}
\begin{prop}\ollabel{ranksexist}
	$\setrank{A}$ existe, para qualquer $A$.
\end{prop}
\begin{proof}
	Fica como exercício.
\end{proof}
\begin{prob}
	Prove \olref[sth][spine][rank]{ranksexist}.	
\end{prob}
\noindent 
A boa ordenação dos postos permite provar alguns resultados importantes:
\begin{prop}\ollabel{valphalowerrank}
Para qualquer ordinal $\alpha$, $V_\alpha = \Setabs{x}{\setrank{x} \in \alpha}$.
\end{prop}

\begin{proof}
Se $\setrank{x} \in \alpha$ então $x \subseteq V_{\setrank{x}} \in
V_\alpha$, logo $x \in V_\alpha$ pois $V_\alpha$ é potente (invocando
\olref[Valphabasic]{Valphabasicprops} várias vezes). Reciprocamente, se $x \in V_\alpha$ então $x \subseteq V_\alpha$, logo $\setrank{x} \leq \alpha$; agora uma indução transfinita simples mostra que $x \notin V_\alpha$. 
\end{proof}
\begin{prob}
	Complete a indução transfinita simples mencionada na \olref[sth][spine][rank]{valphalowerrank}.
\end{prob}

\begin{prop}\ollabel{rankmemberslower}
Se $B \in A$, então $\setrank{B} \in \setrank{A}$.
\end{prop}

\begin{proof}
$A \subseteq V_{\setrank{A}} = \Setabs{x}{\setrank{x} \in
\setrank{A}}$ por \olref{valphalowerrank}.
\end{proof}
\noindent
Com este fato, podemos estabelecer um resultado que nos permite provar
afirmações sobre \emph{todos os conjuntos} por uma forma de indução:

\begin{thm}[Esquema de indução-$\in$] 
Para qualquer fórmula $\phi$:
\[
	\forall A((\forall x \in A)\phi(x) \lif \phi(A)) \lif \forall A \phi(A).
\]
\end{thm}

\begin{proof}
Provaremos a contrapositiva. Suponhamos, pois, $\lnot \forall A
\phi(A)$. Pela indução transfinita
(\olref[ordinals][basic]{ordinductionschema}), existe
algum não-$\phi$ de posto mínimo possível; isto é, algum $A$
tal que $\lnot \phi(A)$ e $\forall x(\setrank{x} \in \setrank{A}
\lif \phi(x))$. Ora, se $x \in A$ então $\setrank{x} \in
\setrank{A}$, por \olref{rankmemberslower}, logo $\phi(x)$; isto é, $(\forall x \in
A)\phi(x) \land \lnot \phi(A)$.
\end{proof}\noindent Eis uma forma informal de interpretar este resultado
poderoso. Digamos que $\phi$ é \emph{hereditária} sse sempre que todo o
!!{element} de um conjunto é $\phi$, o próprio conjunto é $\phi$. Então
a indução-$\in$ diz o seguinte: se $\phi$ é hereditária,
todo o conjunto é $\phi$.

Para concluir a discussão dos postos (por agora), provaremos algumas
afirmações que já antecipamos várias vezes. 

\begin{prop}\ollabel{ranksupstrict}
$\setrank{A} = \supstrict_{x \in A}\setrank{x}$.
\end{prop}

\begin{proof}
Seja $\alpha = \supstrict_{x \in A}\setrank{x}$. Por
\olref{rankmemberslower}, $\alpha \leq \setrank{A}$. Mas se $x \in A$
então $\setrank{x} \in \alpha$, logo $x \in V_\alpha$ por \olref{valphalowerrank}, e portanto $A
\subseteq V_\alpha$, isto é, $\setrank{A} \leq \alpha$. Logo
$\setrank{A} = \alpha$.
\end{proof}

\begin{cor}\ollabel{ordsetrankalpha}
Para qualquer ordinal $\alpha$, $\setrank{\alpha} = \alpha$.
\end{cor}

\begin{proof}
Suponhamos, por indução transfinita, que $\setrank{\beta} = \beta$ para
todo o $\beta \in \alpha$. Ora, $\setrank{\alpha} = \supstrict_{\beta \in
\alpha}\setrank{\beta} = \supstrict_{\beta \in \alpha}\beta = \alpha$
por \olref{ranksupstrict}.
%	Primeiro, notemos que $\setrank{\alpha} \neq \beta$ para qualquer $\beta \in
%	\alpha$. Caso contrário teríamos $\beta = \setrank{\beta} \in
%	\setrank{\alpha} = \beta$ por \olref{rankmemberslower}, isto é,
%	$\beta \in \beta$, uma contradição. 
%
%	Notemos agora que $\alpha \subseteq V_\alpha$. Pois $\alpha =
%	\Setabs{\beta}{\beta \in \alpha}$ por
%	\olref[sfr][ordinals][basic]{ordissetofsmallerord}. E se
%	$\beta \in \alpha$ então $\beta \subseteq V_{\setrank{\beta}} =
%	V_\beta \in V_\alpha$, logo $\beta \in V_\alpha$, por
%	\olref[sfr][spine][Valphabasic]{Valphabasicprops} duas vezes. 
\end{proof}

Finalmente, eis uma prova rápida do resultado prometido no fim de
\olref[foundation]{sec}, a saber que $\ZFminus$ prova a implicação
$\emph{Regularidade} \Rightarrow \emph{Fundação}$. (Notemos que a
noção de «posto» e \olref{rankmemberslower} estão disponíveis para uso
nesta prova pois---como indicamos no início deste texto---podem ser apresentadas usando $\ZFminus + \text{Regularidade}$.)

\begin{prop}[trabalhando em $\ZFminus + \text{Regularidade}$]\ollabel{zfminusregularityfoundation} A Fundação vale.
\end{prop}

\begin{proof}
Fixemos $A \neq \emptyset$, e algum $B \in A$ de posto mínimo possível. Se
$c \in B$ então $\setrank{c} \in \setrank{B}$ por
\olref{rankmemberslower}, logo $c \notin A$ pela escolha de $B$.
\end{proof}

\end{document}
